% =====================================================================
% Limitations
% =====================================================================
\section{Limitations}
\label{sec:limitations}

\paragraph{Scope of the transfer.}
The method assumes that the adapted modules have the same input and output
dimensions in $M_v$ and $M_{v+1}$ and that the two checkpoints share a
tokenizer, or at least tokenize the anchor text identically
(\S\ref{sec:assumptions}; we check the latter mechanically before every run).
Both hold for the version pairs we study and for most minor releases within a
family, but they exclude cross-size transfer (e.g.\ 1.5B\,$\to$\,3B), cross-family
transfer, and releases that change the vocabulary. Ridge regression itself
accepts rectangular maps, so a dimension change is not a mathematical obstacle,
but we have not evaluated it and make no claim there. Cross-LoRA handles the
dimension-mismatch case by construction and remains the method of choice for
heterogeneous pairs; our claim is confined to version drift.

\paragraph{Linear and per-module.}
The transplant is linear so that it folds into the LoRA factors without
changing rank or inference cost. Drift that is nonlinear, or that redistributes
a function across modules (a head that moves to a different layer), cannot be
captured by any per-module linear map, aligned or not. The controlled
experiments in \S\ref{sec:ablations} plant linear drift and so measure the
best case for the method; the real-pair results are the honest measurement of
how much of actual drift is linear enough to be recovered.

\paragraph{Breadth of evidence.}
%TODO(GPU): revise once the 1.5B (and, if obtained, Gemma 3->4) results exist.
The real-model results come from one family and a small number of version
pairs and tasks. The drift study at 0.5B uses anchor sets of a few thousand
tokens on a CPU; the 1.5B experiments use larger anchor sets but still a
single generic corpus. We have not tested whether the anchor domain matters
(code versus prose, or in-domain versus out-of-domain anchors), whether the
conclusions hold for adapters on the MLP projections, or for other PEFT
families (DoRA, IA$^3$, prefix tuning). Larger version jumps than the ones
studied here, for instance a release that changes the attention layout or adds
layers, break the module correspondence the method relies on.

\paragraph{Evaluation.}
Perplexity on held-out task text is a weak proxy for the utility of an adapter,
and on tasks where fine-tuning moves perplexity only slightly the oracle-naive
gap is small, making the gap-recovered metric noisy. We therefore emphasize
tasks with a large oracle lift and an exact-match metric (GSM8K), but
exact-match on greedy decodes with a fixed prompt format is itself sensitive to
formatting changes between versions, which could favour or penalize any
transferred adapter for reasons unrelated to the alignment. Human or
LLM-judged generative quality is not evaluated.

\paragraph{Oracle definition.}
The oracle is a LoRA of the same rank trained on $M_{v+1}$ with the same
recipe and step budget as the source adapter. A better-tuned oracle would
widen the gap and lower every method's recovered fraction; a worse one would
do the opposite. The oracle is a reference point for a fixed training budget,
not a ceiling on what task data can buy.

\paragraph{Cost of anchors.}
Although anchors are unlabeled and generic, the method still needs one forward
pass of the anchor set through both base models with hooks on every adapted
module, and it stores activations of size $N\times d$ per module. For a 12B
model with hundreds of adapted modules this is tens of gigabytes of cached
activations at $N=10^4$; streaming the covariance statistics instead of the
raw activations removes the storage cost but we have not implemented it.
